\documentclass[11pt]{article}
\usepackage[margin=1in]{geometry}
\usepackage{amsmath,amssymb,amsfonts}
\usepackage{booktabs}
\usepackage{hyperref}

\title{QCD Running of the $b$-Quark Yukawa Coupling\\in DFD Model B}
\author{DFD Cross-Reference Analysis}
\date{March 2026}

\begin{document}
\maketitle

\section{Setup}

Model~B assigns a geometric exponent $n_b = 1$ to the $b$-quark
(from $k_f = 1$, $k_H = +1$), giving the bare mass formula
\begin{equation}
  m_b^{\text{bare}} = A_b \,\alpha \,\frac{v}{\sqrt{2}}
  = A_b \times 1270.45 \;\text{MeV},
\end{equation}
where $\alpha = 1/137.036$ and $v/\sqrt{2} = 174\,097.5$~MeV.
The observed $\overline{\text{MS}}$ mass is $m_b(m_b) = 4180$~MeV,
requiring
\begin{equation}
  A_b = \frac{4180}{1270.45} \approx 3.290.
\end{equation}
The reconciliation agent suggested $A_b^{(B)} \approx \pi$.
This note computes how much of this prefactor can be attributed to
QCD running of the Yukawa coupling.

\section{QCD Mass Running}

The $\overline{\text{MS}}$ running quark mass obeys
\begin{equation}
  \frac{m(\mu_1)}{m(\mu_2)}
  = \left[\frac{\alpha_s(\mu_1)}{\alpha_s(\mu_2)}\right]^{d_m(n_f)},
  \qquad
  d_m(n_f) = \frac{12}{33 - 2n_f},
\end{equation}
at one-loop, where the mass anomalous dimension coefficient is
$\gamma_0 = 8$ (for $C_F = 4/3$).  We run $\alpha_s$ using the
two-loop beta function with flavor thresholds at
$m_c = 1.27$~GeV, $m_b = 4.18$~GeV, $m_t = 172.5$~GeV, starting from
$\alpha_s(M_Z) = 0.1179$.

\subsection{$\alpha_s$ Profile}

\begin{center}
\begin{tabular}{rcc}
\toprule
$\mu$ (GeV) & $\alpha_s$ (2-loop) & $n_f$ \\
\midrule
1.27   & 0.3428 & 4 \\
2.00   & 0.2805 & 4 \\
4.18   & 0.2173 & 4 \\
10.0   & 0.1752 & 5 \\
91.2   & 0.1179 & 5 \\
174.1  & 0.1077 & 6 \\
\bottomrule
\end{tabular}
\end{center}

Note: Our 2-loop $\alpha_s(m_b) = 0.217$ is slightly below the PDG
central value of $\sim 0.227$, and $\alpha_s(m_c) = 0.343$ is below
$\sim 0.39$.  A more refined 3-loop treatment would increase both,
amplifying the running factor modestly.

\subsection{Running Factor: $v/\sqrt{2} \to m_b$}

The DFD Yukawa is defined at $\mu_{\text{high}} = v/\sqrt{2} = 174.1$~GeV.
Running down to $\mu = m_b = 4.18$~GeV crosses two segments:

\begin{center}
\begin{tabular}{lccc}
\toprule
Segment & $n_f$ & $d_m$ & Running factor (2-loop) \\
\midrule
$174.1 \to 172.5$~GeV & 6 & 0.5714 & 1.0006 \\
$172.5 \to 4.18$~GeV  & 5 & 0.5217 & 1.4417 \\
\midrule
\textbf{Total} & & & \textbf{1.4426} \\
\bottomrule
\end{tabular}
\end{center}

Including the segment from $m_b$ to $m_c = 1.27$~GeV ($n_f = 4$,
$d_m = 0.48$) adds a further factor 1.245, giving a total running
factor of $R = 1.796$ from 174~GeV to 1.27~GeV.

\section{Results for Model B}

\subsection{With $A_b = 1$ (Pure Running)}

If the DFD formula at the high scale gives
$m_b(\mu_{\text{high}}) = \alpha \cdot v/\sqrt{2} = 1270$~MeV,
then after QCD running:
\begin{equation}
  m_b(m_b) = 1.443 \times 1270 = 1833 \;\text{MeV},
\end{equation}
far short of the observed 4180~MeV.  QCD running alone accounts for
a factor of $\sim$1.44, but the needed prefactor is 3.29.

\subsection{Required Prefactor with Running}

\begin{equation}
  A_b^{\text{needed}} = \frac{m_b(m_b)}{R_{m_b} \times \alpha \cdot v/\sqrt{2}}
  = \frac{4180}{1.443 \times 1270.5}
  \approx 2.281.
\end{equation}

\subsection{Comparison with $A_b = \pi$}

With $A_b = \pi$:
\begin{equation}
  m_b(m_b)\big|_{A_b=\pi} = 1.443 \times \pi \times 1270.5
  = 5758 \;\text{MeV},
\end{equation}
which overshoots the observed value by 37.7\%.
Thus $A_b = \pi$ is \textbf{not} consistent with Model~B when
QCD running is included.

\section{Model A Comparison}

Model~A uses $n_b = 0$ and $A_b = 1/42$:
\begin{equation}
  m_b = \frac{1}{42}\,\frac{v}{\sqrt{2}} = 4145 \;\text{MeV},
\end{equation}
matching the observed 4180~MeV to 0.8\%, with no QCD running
correction needed (the formula is already interpreted as a pole or
low-scale mass).

\section{Conclusions}

\begin{enumerate}
\item QCD running from $v/\sqrt{2} = 174$~GeV down to $m_b = 4.18$~GeV
      provides a mass enhancement factor $R \approx 1.44$ (2-loop).
\item This is substantial but insufficient to bridge the gap from
      $A_b = 1$ to the needed $A_b \approx 3.29$.  The running
      reduces the required prefactor to $A_b \approx 2.28$.
\item The suggestion $A_b = \pi \approx 3.14$ overshoots by $\sim$38\%
      when combined with QCD running.  The value $A_b = \pi$ would
      work \emph{without} QCD running (giving 3991~MeV, off by 4.5\%),
      but this ignores mandatory QCD corrections.
\item A 3-loop treatment with proper matching would increase $R$
      modestly (perhaps to $\sim$1.55--1.6), lowering $A_b^{\text{needed}}$
      to $\sim$2.1, still not matching $\pi$.
\item Model~A's direct assignment $m_b = v/(42\sqrt{2})$ remains
      significantly more accurate (0.8\% vs Model~B's need for an
      unexplained $O(1)$ prefactor).
\item If Model~B is pursued, the required prefactor $A_b \approx 2.28$
      after QCD running needs a group-theoretic or representation-theoretic
      origin distinct from $\pi$.
\end{enumerate}

\end{document}
